% s6_conclusion.tex
% Status: DRAFT — generated by TAR-Author; review before locking
% Lock condition: after Phase 10 numbers finalised

\section{Conclusion}
\label{sec:conclusion}

We have introduced Thermodynamic Continual Learning (TCL), a method that uses per-task
activation entropy as a real-time consolidation signal to govern weight anchoring in
sequential learning.
The core mechanism --- anchoring the task's activation standard deviation at the start
of training, then detecting the ordered regime when activations fall below that anchor
--- is lightweight, requires no Fisher matrix computation, and provides a per-step
control signal rather than a retrospective importance estimate.

On Split-CIFAR-10 with a ResNet-18 backbone across five random seeds, TCL achieves
mean forgetting $0.1161 \pm 0.029$ against EWC's $0.1909 \pm 0.023$
(mean delta $-0.0748$, $p = 0.0190$, Cohen's $d = 1.70$, 5/5 seeds; controlled-rerun
Outcome~A) and against SGD's $0.2331 \pm 0.003$
($p = 0.0012$, $d = 3.68$, 5/5 seeds).
A pre-registered four-condition ablation shows that the anchor penalty is the
load-bearing component: governor-only tracks SGD, penalty-only captures most of the
forgetting reduction, and full TCL remains best on mean forgetting, though the
full-vs.-penalty gap is not yet resolved at $n=5$.

We have also demonstrated a failure mode of forgetting-only evaluation: Synaptic
Intelligence at published default hyperparameters collapses to chance-level accuracy
($0.500$) on all five seeds while reporting competitive forgetting scores.
The joint accuracy-forgetting (JAF) metric proposed here exposes this failure with
minimal reporting overhead and no new hyperparameters.

The thermodynamic framing of neural network consolidation opens a broader research
direction: activation entropy provides a task-agnostic, parameter-group-level signal
that does not require importance score storage, scales linearly in model size, and
can in principle operate in class-incremental and domain-incremental settings without
architectural change.
We view the Phase~9 and Phase~10 results as proof-of-concept at realistic scale;
extending the evaluation to harder settings and longer task sequences is the natural
next step.
